\subsection{The N-Line Circuit Equations}
% ------------------------------------------------------------

Refer back to the general model of \figref{fig:P1Fig1}.
For line~$n$, KVL from the left input terminal through the series
$R$--$L$ and the shunt capacitor to ground gives:
\[
    v_{nP} + v_{Pg}
    = \Bigl(R + pL + \frac{S_n}{p}\Bigr)\,i_n,
    \qquad S_n = \frac{1}{C_n},
\]
where $p = d/dt$ is the Heaviside operator and $S_n$ is the
\emph{elastance} (reciprocal capacitance) of line $n$'s shunt element.

\begin{physbox}[Why Elastance?]
Writing $S_n/p$ instead of $1/(pC_n)$ separates the circuit
parameter ($S_n$) from the integrating operator ($1/p$), making
the matrix algebra cleaner.  The elastance matrix $\mathbf{S}_N$
collects all $N$ shunt-cap parameters into a single object.
\end{physbox}

Stacking $N$ such equations into column-vector form:

\begin{eqnbox}[Mixed-Mode Circuit Equations — Eq.~9 in paper]
\begin{equation}
    \mathbf{Z}\,\mathbf{i}
    = \mathbf{v}_{nP} + v_{Pg}\,\mathbf{1}_N
    \label{eq:MMeqn}
\end{equation}
where the impedance operator matrix and the elastance matrix are
\[
    \mathbf{Z}
    = (R + pL)\,\mathbf{I}_N + \frac{1}{p}\,\mathbf{S}_N,
    \qquad
    \mathbf{S}_N
    = \operatorname{diag}\!\Bigl(
        \underbrace{\tfrac{1}{C},\;\ldots,\;\tfrac{1}{C}}_{N-1},\;
        \tfrac{1}{kC}
      \Bigr),
\]
and $\mathbf{1}_N = [1,\ldots,1]^T$,\;
$\mathbf{i} = [i_1,\ldots,i_N]^T$,\;
$\mathbf{v}_{nP} = [v_{1P},\ldots,v_{NP}]^T$.
\end{eqnbox}

\begin{statebox}[Two Observations About $\mathbf{Z}$]
\begin{itemize}
    \item $(R + pL)\,\mathbf{I}_N$ — every line has the same $R$
          and $L$, so this is a scalar multiple of the identity.
          It will transform cleanly.
    \item $\mathbf{S}_N$ — diagonal, but \emph{not} proportional
          to the identity when $k \neq 1$.  This is the \textbf{only}
          asymmetric object in the equation and the only source of
          DM--CM coupling after transformation.
\end{itemize}
\end{statebox}

% ------------------------------------------------------------
\subsection{Applying the DM--CM Transformation}
% ------------------------------------------------------------

Multiply Eq.~\eqref{eq:MMeqn} on the left by $\TvN$, then
substitute $\mathbf{i} = (\TiN)^{-1}\iDCM$:
\[
    \underbrace{\TvN\,\mathbf{Z}\,(\TiN)^{-1}}_{\mathbf{Z}_{\mathrm{DCM}}}
    \;\iDCM
    \;=\;
    \vDCM
    \;+\; v_{Pg}\,\underbrace{\TvN\,\mathbf{1}_N}_{\mathbf{e}_N}.
\]

\begin{eqnbox}[DCM Equations — Eq.~10 in paper]
\begin{equation}
    \mathbf{Z}_{\mathrm{DCM}}\,\iDCM
    = \vDCM + v_{Pg}\,\mathbf{e}_N
    \label{eq:DCMeqn}
\end{equation}
where $\mathbf{Z}_{\mathrm{DCM}} = \TvN\,\mathbf{Z}\,(\TiN)^{-1}$
and $\mathbf{e}_N = [0,\ldots,0,1]^T$ is the $N$-th unit vector.
\end{eqnbox}

\subsubsection*{Why $\TvN\,\mathbf{1}_N = \mathbf{e}_N$}

Verify for $N = 2$:
\[
\TvN\,\mathbf{1}_2
=
\begin{bmatrix}
1 & -1 \\ \tfrac{1}{2} & \tfrac{1}{2}
\end{bmatrix}
\begin{bmatrix}1 \\ 1\end{bmatrix}
=
\begin{bmatrix}1-1 \\ \tfrac{1}{2}+\tfrac{1}{2}\end{bmatrix}
=
\begin{bmatrix}0 \\ 1\end{bmatrix}
= \mathbf{e}_2.
\]
Every DM row of $\TvN$ has a $+1$ and a $-1$ in adjacent columns
(difference of a constant = 0); the CM row has all $+1/N$ (sum = 1).

\begin{statebox}[$v_{Pg}$ Lives Only in the CM Equation]
Because $v_{Pg}$ is multiplied by $\mathbf{e}_N$, it appears in the
\textbf{CM equation only} — absent from every DM equation.
DM quantities are completely independent of the ground connection.
This confirms the physical intuition: DM current is a loop current
confined to the wire bundle; only the CM path connects to ground
through $v_{Pg}$.
\end{statebox}

\subsubsection*{Expanding $\mathbf{Z}_{\mathrm{DCM}}$}

Since $\mathbf{Z} = (R + pL)\,\mathbf{I}_N + (1/p)\,\mathbf{S}_N$,
linearity gives:
\[
    \mathbf{Z}_{\mathrm{DCM}}
    = (R + pL)\underbrace{\TvN(\TiN)^{-1}}_{\displaystyle\mathbf{M}}
    + \frac{1}{p}\underbrace{\TvN\,\mathbf{S}_N\,(\TiN)^{-1}}
                             {\displaystyle\mathbf{S}_{\mathrm{DCM}}}.
\]

\paragraph{$\mathbf{M} = \TvN(\TiN)^{-1}$ is always diagonal.}

Computing for $N = 2$:
\[
    (\TiN)^{-1}
    = \begin{bmatrix}1 & \tfrac{1}{2} \\[3pt] -1 & \tfrac{1}{2}\end{bmatrix},
    \qquad
    \mathbf{M}
    =
    \begin{bmatrix}1 & -1 \\[3pt] \tfrac{1}{2} & \tfrac{1}{2}\end{bmatrix}
    \begin{bmatrix}1 & \tfrac{1}{2} \\[3pt] -1 & \tfrac{1}{2}\end{bmatrix}
    =
    \begin{bmatrix}2 & 0 \\[3pt] 0 & \tfrac{1}{2}\end{bmatrix}.
\]

The entries are $\Ni = 2$ for the DM row and $1/N = 1/2$ for the CM
row.  For all $N$:
$\mathbf{M} = \operatorname{diag}(\Ni,\ldots,\Ni,\,1/N)$.

\begin{circuitbox}[Physical Meaning of $\mathbf{M}$]
\begin{itemize}
    \item DM: $R_{\mathrm{DM}} = \Ni R = 2R$,\; $L_{\mathrm{DM}} = 2L$
          — DM current flows through both $R$'s in series. \checkmark
    \item CM: $R_{\mathrm{CM}} = R/N = R/2$,\; $L_{\mathrm{CM}} = L/2$
          — CM current splits equally across both lines in parallel. \checkmark
    \item \textbf{No off-diagonal entries.}  $R$ and $L$ never couple
          DM to CM regardless of capacitor asymmetry.
\end{itemize}
\end{circuitbox}

\paragraph{$\mathbf{S}_{\mathrm{DCM}} = \TvN\,\mathbf{S}_N\,(\TiN)^{-1}$
depends on $k$.}

This is the key computation.  For $N = 2$ with
$\mathbf{S}_2 = \operatorname{diag}(1/C,\;1/(kC))$:

\textbf{Stage A} — left-multiply by $\TvN$:
\[
\begin{bmatrix}1 & -1 \\[3pt] \tfrac{1}{2} & \tfrac{1}{2}\end{bmatrix}
\begin{bmatrix}\tfrac{1}{C} & 0 \\[3pt] 0 & \tfrac{1}{kC}\end{bmatrix}
=
\begin{bmatrix}
    \tfrac{1}{C}  & -\tfrac{1}{kC} \\[6pt]
    \tfrac{1}{2C} & \tfrac{1}{2kC}
\end{bmatrix}.
\]

\textbf{Stage B} — right-multiply by $(\TiN)^{-1}$:
\[
\begin{bmatrix}
    \tfrac{1}{C}  & -\tfrac{1}{kC} \\[6pt]
    \tfrac{1}{2C} & \tfrac{1}{2kC}
\end{bmatrix}
\begin{bmatrix}1 & \tfrac{1}{2} \\[3pt] -1 & \tfrac{1}{2}\end{bmatrix}
=
\frac{1}{C}
\begin{bmatrix}
    \dfrac{k+1}{k}  & \dfrac{k-1}{2k}  \\[10pt]
    \dfrac{k-1}{2k} & \dfrac{k+1}{4k}
\end{bmatrix}.
\]

\begin{eqnbox}[DCM Elastance Matrix ($N = 2$) — Eq.~16 in paper]
\begin{equation}
\mathbf{S}_{\mathrm{DCM}}
= \frac{1}{C}
  \begin{bmatrix}
      \dfrac{k+1}{k}  & \dfrac{k-1}{2k}  \\[10pt]
      \dfrac{k-1}{2k} & \dfrac{k+1}{4k}
  \end{bmatrix}
\label{eq:Sdcm}
\end{equation}
Off-diagonal entry $\tfrac{1}{C}\cdot\tfrac{k-1}{2k}$ is the
\textbf{DM--CM coupling elastance}.
It is zero if and only if $k = 1$.
\end{eqnbox}

% ------------------------------------------------------------
\subsection{Case 1 — Symmetric ($k = 1$): Modes Fully Decouple}
% ------------------------------------------------------------

Substituting $k = 1$ into Eq.~\eqref{eq:Sdcm}:
\[
\mathbf{S}_{\mathrm{DCM}}\big|_{k=1}
=
\frac{1}{C}\begin{bmatrix}2 & 0 \\ 0 & \tfrac{1}{2}\end{bmatrix}.
\]
Off-diagonal = 0.  The two rows of Eq.~\eqref{eq:DCMeqn} are
completely independent:

\begin{eqnbox}[Symmetric Decoupled Equations — Eqs.~14 and 15 in paper]
\begin{align}
    \Bigl(2R + 2pL + \frac{2}{pC}\Bigr)\,i_{12}
        &= v_{12}
        \label{eq:DM_sym} \\[6pt]
    \Bigl(\frac{R}{2} + \frac{pL}{2} + \frac{1}{2pC}\Bigr)\,\icm
        &= \vcm + v_{Pg}
        \label{eq:CM_sym}
\end{align}
\end{eqnbox}

\begin{circuitbox}[Reading the Symmetric Equations]
\begin{itemize}
    \item \textbf{DM (Eq.~\ref{eq:DM_sym}):} $i_{12}$ sees
          \emph{series} elements — $2R$, $2L$, and elastance $2/C$
          (two capacitors $C$ in series give $C/2$, elastance $2/C$).
          No $v_{Pg}$ — DM is ground-independent.
    \item \textbf{CM (Eq.~\ref{eq:CM_sym}):} $\icm$ sees
          \emph{parallel} elements — $R/2$, $L/2$, and elastance
          $1/(2C)$ (two $C$'s in parallel give $2C$, elastance $1/(2C)$).
          $v_{Pg}$ drives the CM current — this is the EMI path.
    \item \textbf{No coupling:} a DM voltage source excites zero CM
          current, and vice versa.  This is the ideal symmetric world.
\end{itemize}
\end{circuitbox}

% ------------------------------------------------------------
\subsection{Case 2 — Asymmetric ($k \neq 1$): Coupling Appears}
% ------------------------------------------------------------

With $k \neq 1$, the off-diagonal entries of $\mathbf{S}_{\mathrm{DCM}}$
are non-zero.  Expanding Eq.~\eqref{eq:DCMeqn} row by row:

\begin{eqnbox}[Asymmetric Coupled Equations ($N = 2$)]
\begin{align}
    \Bigl(2R + 2pL + \frac{k+1}{kpC}\Bigr)\,i_{12}
    + \frac{k-1}{2kpC}\,\icm
        &= v_{12}
        \label{eq:DM_asym} \\[6pt]
    \frac{k-1}{2kpC}\,i_{12}
    + \Bigl(\frac{R}{2} + \frac{pL}{2} + \frac{k+1}{4kpC}\Bigr)\,\icm
        &= \vcm + v_{Pg}
        \label{eq:CM_asym}
\end{align}
\end{eqnbox}

The term $\tfrac{k-1}{2kpC}$ appears in \emph{both} equations —
it links $i_{12}$ and $\icm$ bidirectionally.

\begin{warningbox}[The Coupling Term $\dfrac{k-1}{2kpC}$]
\begin{itemize}
    \item $k = 1$: coupling $= 0$.  Eqs.~\eqref{eq:DM_asym} and
          \eqref{eq:CM_asym} reduce exactly to
          Eqs.~\eqref{eq:DM_sym} and \eqref{eq:CM_sym}. \checkmark
    \item $k < 1$ (real MCPM — source cap is smaller):
          coupling is negative.  A positive DM oscillation drives
          CM current in the opposite direction.
    \item The coupling is \emph{purely capacitive} (the $1/p$ factor
          means it behaves like a capacitor element in the circuit).
          Changing $R$ or $L$ cannot remove it — only restoring
          $k = 1$ sets the coupling to zero.
\end{itemize}
\end{warningbox}

% ------------------------------------------------------------
\subsection{Numerical Proof — The 2\texorpdfstring{$\times$}{×}2
            Case Study}
% ------------------------------------------------------------

The algebra of \S2.1--\S2.4 shows \emph{why} the coupling term exists.
The most direct verification is a concrete number: start from
Eq.~\eqref{eq:Sdcm}, pick real component values, and check whether
a DM excitation alone can produce a CM current.

\medskip
\textbf{The question we answer:}
\begin{center}
\emph{Apply a 1\,V sinusoidal DM voltage with no CM source.
      How much CM current flows?}
\end{center}
If the answer is zero the modes are decoupled; if non-zero, then a
DM switching waveform is generating conducted EMI by itself — with
no CM voltage source required.

\subsubsection*{Setup (using \S1 notation throughout)}

Recall from \S1.5 that the $N = 2$ decomposition encodes:
$i_{12} = (i_1 - i_2)/2$ (DM) and $\icm = i_1 + i_2$ (CM),
with the same split for voltages.
We now work entirely in this DM--CM coordinate system.

Choose the simplest parameter set that isolates the capacitor effect:

\begin{itemize}
    \item \textbf{Shunt capacitances:} $C = 1\,\text{nF}$ for
          line~1; $kC$ for line~2 (the asymmetric cap introduced
          in \figref{fig:P1Fig1}).
    \item \textbf{DM excitation:} $V_{12} = 1\angle 0°\,\text{V}$
          at $f = 1\,\text{MHz}$ ($\omega = 2\pi f \approx
          6.28 \times 10^6\,\text{rad/s}$).  Phasor notation:
          $p \to j\omega$.
    \item \textbf{No CM source:} $V_{\mathrm{CM}} = 0$ and
          $V_{Pg} = 0$.  This is the critical choice — we want
          to see whether DM alone drives $\icm$.
    \item \textbf{Series elements:} $R = L = 0$ to isolate the
          capacitor contribution.  (Including $R$ and $L$ only
          adds the diagonal matrix $\mathbf{M}$ from \S2.2 —
          diagonal matrices never couple DM to CM, so they
          cannot change the conclusion.)
\end{itemize}

With $R = L = 0$ and $p \to j\omega$, Eq.~\eqref{eq:DCMeqn}
reduces to
$\mathbf{S}_{\mathrm{DCM}}/(j\omega) \cdot \mathbf{I}_{\mathrm{DCM}}
= \mathbf{V}_{\mathrm{DCM}}$.
Substituting Eq.~\eqref{eq:Sdcm} gives the phasor system we will
solve twice — once for $k = 1$ and once for $k = 0.5$:

\begin{eqnbox}[2-Line Phasor System Under Pure DM Excitation]
\begin{equation}
\frac{1}{j\omega C}
\begin{bmatrix}
    \dfrac{k+1}{k}  & \dfrac{k-1}{2k} \\[10pt]
    \dfrac{k-1}{2k} & \dfrac{k+1}{4k}
\end{bmatrix}
\begin{bmatrix} I_{12} \\[4pt] I_{\mathrm{CM}} \end{bmatrix}
=
\begin{bmatrix} 1\,\mathrm{V} \\[4pt] 0 \end{bmatrix}
\label{eq:num_sys}
\end{equation}
Right-hand side: $[V_{12},\;V_{\mathrm{CM}}+V_{Pg}]^T = [1\,\mathrm{V},\;0]^T$.
A non-zero $I_{\mathrm{CM}}$ in the solution means the DM voltage
alone is generating CM current — the definition of asymmetry-induced
EMI.
\end{eqnbox}

% ---- Case A ----
\subsubsection*{Case A — Symmetric ($k = 1$): modes stay silent}

Substitute $k = 1$.  The off-diagonal of Eq.~\eqref{eq:Sdcm}
becomes $\tfrac{k-1}{2k}\big|_{k=1} = 0$, so
Eq.~\eqref{eq:num_sys} becomes:
\[
\frac{1}{j\omega C}
\begin{bmatrix} 2 & 0 \\ 0 & \tfrac{1}{2} \end{bmatrix}
\begin{bmatrix} I_{12} \\ I_{\mathrm{CM}} \end{bmatrix}
=
\begin{bmatrix} 1\,\mathrm{V} \\ 0 \end{bmatrix}.
\]
The matrix is block-diagonal: \textbf{Row 1 and Row 2 are
completely independent.}  Row 2 alone gives:
\[
    \frac{1}{2j\omega C}\,I_{\mathrm{CM}} = 0
    \quad\Longrightarrow\quad
    \boxed{I_{\mathrm{CM}} = 0.}
\]
This holds for \emph{any} value of $V_{12}$.  The DM source is
completely invisible to the CM equation.  Row 1 gives
$I_{12} = j\omega C \cdot V_{12}/2$ (a nonzero DM current that
circulates inside the wire bundle), but it drives zero ground
current — exactly Eq.~\eqref{eq:DM_sym}.

\begin{statebox}[What Decoupling Means in Practice]
With $k = 1$, no matter how large or fast the DM switching
transient is, $\icm = 0$ and no asymmetry-induced EMI reaches the
LISN.  The only remaining EMI path is $V_{Pg}$ on the right-hand
side of Eq.~\eqref{eq:CM_sym} — driven by the CM voltage of the
power source, not by the DM switching waveform.
\end{statebox}

% ---- Case B ----
\subsubsection*{Case B — Asymmetric ($k = 0.5$): DM drives CM current}

SiC power-module source terminals typically have roughly half the
parasitic capacitance of the drain terminal (see \S0.5), making
$k \approx 0.5$ a realistic value.
Substitute $k = 0.5$ into Eq.~\eqref{eq:Sdcm}:

\begin{center}
\renewcommand{\arraystretch}{1.5}
\begin{tabular}{lcc}
\toprule
Matrix entry & Formula & Value at $k = 0.5$ \\
\midrule
DM self-elastance  & $(k+1)/k$    & $1.5/0.5 = 3$ \\
Coupling (off-diag) & $(k-1)/(2k)$ & $-0.5/1.0 = -\,0.5$ \\
CM self-elastance  & $(k+1)/(4k)$ & $1.5/2.0 = 0.75$ \\
\bottomrule
\end{tabular}
\end{center}

The phasor system (Eq.~\ref{eq:num_sys}) becomes:
\[
\frac{1}{j\omega C}
\begin{bmatrix} 3 & -0.5 \\ -0.5 & 0.75 \end{bmatrix}
\begin{bmatrix} I_{12} \\ I_{\mathrm{CM}} \end{bmatrix}
=
\begin{bmatrix} 1\,\mathrm{V} \\ 0 \end{bmatrix}.
\]

\textbf{Step 1 — solve Row 2 for $I_{\mathrm{CM}}$:}
\[
    -0.5\,I_{12} + 0.75\,I_{\mathrm{CM}} = 0
    \quad\Longrightarrow\quad
    I_{\mathrm{CM}} = \frac{0.5}{0.75}\,I_{12} = \frac{2}{3}\,I_{12}.
\]
The CM and DM currents are locked together by the off-diagonal term.

\textbf{Step 2 — substitute into Row 1:}
\[
    3\,I_{12} + (-0.5)\cdot\tfrac{2}{3}\,I_{12}
    = j\omega C \cdot 1\,\mathrm{V}
    \;\;\Longrightarrow\;\;
    \Bigl(3 - \tfrac{1}{3}\Bigr)I_{12} = j\omega C
    \;\;\Longrightarrow\;\;
    \frac{8}{3}\,I_{12} = j\omega C.
\]
\[
    I_{12} = j\omega C \cdot \frac{3}{8}.
\]

\textbf{Step 3 — back-substitute to get $I_{\mathrm{CM}}$:}
\[
    I_{\mathrm{CM}}
    = \frac{2}{3}\,I_{12}
    = \frac{2}{3} \cdot j\omega C \cdot \frac{3}{8}
    = \frac{j\omega C}{4}.
\]

\begin{eqnbox}[CM Current Generated by a Pure DM Source ($k=0.5$, $N=2$)]
\begin{equation}
    I_{\mathrm{CM}} = \frac{j\omega C}{4}\cdot V_{12}
    \label{eq:ICM_result}
\end{equation}
\end{eqnbox}

\textbf{Plug in numbers:} $C = 1\,\text{nF}$, $f = 1\,\text{MHz}$,
$V_{12} = 1\,\text{V}$:
\[
|I_{\mathrm{CM}}|
= \frac{2\pi \times 10^6 \times 10^{-9}}{4}
= \frac{6.28 \times 10^{-3}}{4}
\approx \mathbf{1.57\,\text{mA}}.
\]
A 1\,V DM voltage at 1\,MHz with $k = 0.5$ produces
\textbf{1.57\,mA of CM current} — with no CM source, no $V_{Pg}$,
purely from the capacitor mismatch.

\begin{warningbox}[Connecting Back to \S1: What Does 1.57\,mA Mean?]
From \S1.6 (Eq.~\ref{eq:inv_i1}) each line carries $\icm/2$ as its
CM share.  So $I_{\mathrm{CM}} = 1.57\,\text{mA}$ means line~1
carries $0.785\,\text{mA}$ more than line~2.  That imbalance is the
\emph{same current} that the $\icm$ definition (Eq.~\ref{eq:icm})
would measure with a clamp-meter around both wires — exactly the
conducted EMI the LISN captures (recall \S0.4).
The DM switching waveform acts as the pump; the capacitor
mismatch ($k \neq 1$) is the valve that lets a fraction of it
escape to the ground return.
\end{warningbox}

\begin{statebox}[Side-by-Side Summary: $k = 1$ vs.\ $k = 0.5$]
\begin{center}
\renewcommand{\arraystretch}{1.5}
\begin{tabular}{lcc}
\toprule
\textbf{Quantity}
    & $k = 1$ (symmetric)
    & $k = 0.5$ (asymmetric) \\
\midrule
Off-diag of $\mathbf{S}_{\mathrm{DCM}}$
    & $0$ & $-0.5/C$ \\
Equations coupled?
    & No  & Yes \\
$I_{12}$ (DM current)
    & $j\omega C/2$
    & $j\omega C \cdot 3/8$ \\
$I_{\mathrm{CM}}$ given $V_{12}=1\,\text{V}$, $V_{\mathrm{CM}}=0$
    & $0$
    & $j\omega C/4$ \\
$|I_{\mathrm{CM}}|$ at $C=1\,\text{nF}$, $f = 1\,\text{MHz}$
    & $0\,\text{mA}$
    & $1.57\,\text{mA}$ \\
Asymmetry-induced EMI?
    & \textbf{None}
    & \textbf{Yes — from DM alone} \\
\bottomrule
\end{tabular}
\end{center}
Same DM excitation, identical $R$, $L$, $C$ values — only the
ratio $k$ changed.  The CM current jumps from zero to 1.57\,mA.
This is the numerical proof of the opening paragraph's claim.
\end{statebox}

% ---- N = 3 extension ----
\subsubsection*{Extension to $N = 3$ Lines}

The $N = 2$ proof captures the core result.  Extending to $N = 3$
reveals an additional effect that $N = 2$ cannot show: when one
capacitor is $kC$, asymmetry couples not only DM to CM, but also
the two DM pairs \emph{to each other}.

\subsubsection*{Step 1 — The $N = 3$ DCM Elastance Matrix}

With $\mathbf{S}_3 = \operatorname{diag}(1/C,\;1/C,\;1/(kC))$
(lines 1 and 2 symmetric, line 3 asymmetric at capacitance $kC$)
and the $N = 3$ matrices from \S1.5, the same two-stage computation
of \S2.2 gives:

\begin{eqnbox}[DCM Elastance Matrix ($N = 3$, general $k$)]
\begin{equation}
\mathbf{S}_{\mathrm{DCM}}^{(3)}
= \frac{1}{C}
\begin{bmatrix}
    2 & 0 & 0 \\[8pt]
    \dfrac{2(1-k)}{3k} & \dfrac{2(k+2)}{3k} & \dfrac{k-1}{3k} \\[12pt]
    \dfrac{2(k-1)}{9k} & \dfrac{4(k-1)}{9k} & \dfrac{2k+1}{9k}
\end{bmatrix}
\label{eq:Sdcm_N3}
\end{equation}
All off-diagonal entries are zero if and only if $k = 1$.
\end{eqnbox}

\begin{statebox}[Structural Observation: Row 1 is Always $k$-Independent]
Row 1 is $[2/C,\;0,\;0]$ for \emph{any} $k$.  Lines 1 and 2 both
have cap $C$, so the DM(1--2) equation never involves $kC$.
Coupling enters only in rows 2 and 3 — DM pair (2--3) and CM —
because both involve line 3.
\end{statebox}

\subsubsection*{Step 2 — The Phasor System}

With $R = L = 0$ and $p \to j\omega$, under pure DM(1--2) excitation
($V_{12} = 1\,\text{V}$, $V_{23} = V_{\mathrm{CM}} = V_{Pg} = 0$):

\begin{eqnbox}[$N = 3$ Phasor System Under DM(1--2) Excitation Only]
\begin{equation}
\frac{1}{j\omega C}
\begin{bmatrix}
    2 & 0 & 0 \\[8pt]
    \dfrac{2(1-k)}{3k} & \dfrac{2(k+2)}{3k} & \dfrac{k-1}{3k} \\[12pt]
    \dfrac{2(k-1)}{9k} & \dfrac{4(k-1)}{9k} & \dfrac{2k+1}{9k}
\end{bmatrix}
\begin{bmatrix} I_{12} \\[4pt] I_{23} \\[4pt] I_{\mathrm{CM}} \end{bmatrix}
=
\begin{bmatrix} 1\,\mathrm{V} \\[4pt] 0 \\[4pt] 0 \end{bmatrix}
\label{eq:num_sys_N3}
\end{equation}
\end{eqnbox}

\subsubsection*{Case A — Symmetric ($k = 1$): three independent equations}

At $k = 1$ every off-diagonal entry vanishes.  The matrix becomes
diagonal and the three rows decouple:
\[
\frac{1}{j\omega C}
\begin{bmatrix} 2 & 0 & 0 \\ 0 & 2 & 0 \\ 0 & 0 & \tfrac{1}{3} \end{bmatrix}
\begin{bmatrix} I_{12} \\ I_{23} \\ I_{\mathrm{CM}} \end{bmatrix}
=
\begin{bmatrix} 1\,\mathrm{V} \\ 0 \\ 0 \end{bmatrix}.
\]
Each row gives one unknown immediately:
\[
    I_{12} = \frac{j\omega C}{2}, \qquad I_{23} = 0, \qquad
    \boxed{I_{\mathrm{CM}} = 0.}
\]
Only the directly excited DM(1--2) current is non-zero.  DM(2--3)
and CM are completely silent — modes are decoupled.

\begin{statebox}[Diagonal Entries at $k = 1$: Circuit Check]
\begin{itemize}
    \item $S_{11} = 2/C$: DM(1--2) cap impedance.
          Lines 1 and 2 each have cap $C$; their DM current sees
          them in series $\Rightarrow$ $C/2$, elastance $= 2/C$. \checkmark
    \item $S_{22} = 2/C$: DM(2--3) cap impedance.
          At $k = 1$ lines 2 and 3 also have cap $C$; same
          series result. \checkmark
    \item $S_{33} = 1/(3C)$: CM cap impedance.
          Three caps $C$ in parallel $\Rightarrow$ $3C$,
          elastance $= 1/(3C)$. \checkmark
\end{itemize}
\end{statebox}

\subsubsection*{Case B — Asymmetric ($k = 0.5$): coupling across all
three equations}

Substitute $k = 0.5$ into Eq.~\eqref{eq:Sdcm_N3}:

\begin{center}
\renewcommand{\arraystretch}{1.5}
\begin{tabular}{lcc}
\toprule
Entry & Formula & Value at $k = 0.5$ \\
\midrule
$S_{21}$ (DM23 driven by DM12) & $2(1-k)/(3k)$  & $2/3$  \\
$S_{22}$ (DM23 self)           & $2(k+2)/(3k)$  & $10/3$ \\
$S_{23}$ (DM23 driven by CM)   & $(k-1)/(3k)$   & $-1/3$ \\
$S_{31}$ (CM driven by DM12)   & $2(k-1)/(9k)$  & $-2/9$ \\
$S_{32}$ (CM driven by DM23)   & $4(k-1)/(9k)$  & $-4/9$ \\
$S_{33}$ (CM self)             & $(2k+1)/(9k)$  & $4/9$  \\
\bottomrule
\end{tabular}
\end{center}

The system (Eq.~\ref{eq:num_sys_N3}) becomes:
\[
\frac{1}{j\omega C}
\begin{bmatrix}
    2 & 0 & 0 \\
    \tfrac{2}{3} & \tfrac{10}{3} & -\tfrac{1}{3} \\
    -\tfrac{2}{9} & -\tfrac{4}{9} & \tfrac{4}{9}
\end{bmatrix}
\begin{bmatrix} I_{12} \\ I_{23} \\ I_{\mathrm{CM}} \end{bmatrix}
=
\begin{bmatrix} 1\,\mathrm{V} \\ 0 \\ 0 \end{bmatrix}.
\]

\textbf{Step 1 — Row 1} (same structure as $k = 1$, since row 1 is
$k$-independent):
\[
    I_{12} = \frac{j\omega C}{2}.
\]

\textbf{Step 2 — substitute $I_{12}$ into Rows 2 and 3}
(multiply each row by $j\omega C$):

Row~2: $\;\tfrac{2}{3}\cdot\tfrac{j\omega C}{2}
        + \tfrac{10}{3}\,I_{23}
        - \tfrac{1}{3}\,I_{\mathrm{CM}} = 0$
\quad$\Rightarrow$\quad multiply by 3:
\[
    j\omega C + 10\,I_{23} - I_{\mathrm{CM}} = 0. \hspace{1.5cm}\text{(A)}
\]

Row~3: $\;-\tfrac{2}{9}\cdot\tfrac{j\omega C}{2}
          - \tfrac{4}{9}\,I_{23}
          + \tfrac{4}{9}\,I_{\mathrm{CM}} = 0$
\quad$\Rightarrow$\quad multiply by 9:
\[
    -j\omega C - 4\,I_{23} + 4\,I_{\mathrm{CM}} = 0. \hspace{1.5cm}\text{(B)}
\]

\textbf{Step 3 — solve (A) and (B) simultaneously:}

From (B): $I_{\mathrm{CM}} = (j\omega C + 4\,I_{23})/4$.
Substitute into (A):
\[
    j\omega C + 10\,I_{23} - \frac{j\omega C + 4\,I_{23}}{4} = 0
    \;\;\Rightarrow\;\;
    \frac{3\,j\omega C}{4} + 9\,I_{23} = 0
    \;\;\Rightarrow\;\;
    I_{23} = -\frac{j\omega C}{12}.
\]
\[
    I_{\mathrm{CM}}
    = \frac{j\omega C + 4(-j\omega C/12)}{4}
    = \frac{j\omega C\,(1 - 1/3)}{4}
    = \frac{j\omega C}{6}.
\]

\begin{eqnbox}[Final Results ($N = 3$, $k = 0.5$, $V_{12} = 1\,\text{V}$,
              $V_{23} = V_{\mathrm{CM}} = 0$)]
\[
    I_{12} = \frac{j\omega C}{2},
    \qquad
    I_{23} = -\frac{j\omega C}{12},
    \qquad
    I_{\mathrm{CM}} = \frac{j\omega C}{6}.
\]
At $C = 1\,\text{nF}$, $f = 1\,\text{MHz}$:\quad
$|I_{12}| \approx 3.14\,\text{mA}$,\;
$|I_{23}| \approx 0.52\,\text{mA}$,\;
$|I_{\mathrm{CM}}| \approx 1.05\,\text{mA}$.
\end{eqnbox}

\begin{physbox}[Physical Check: Why $I_{\mathrm{CM}} \neq 0$ for $k = 0.5$]
The excitation $V_{12} = 1\,\text{V}$, $V_{23} = V_{\mathrm{CM}} = 0$
forces the three line voltages to:
\[
    v_{1P} = +\tfrac{2}{3}\,\text{V},\qquad
    v_{2P} = -\tfrac{1}{3}\,\text{V},\qquad
    v_{3P} = -\tfrac{1}{3}\,\text{V}.
\]
With $v_{Pg} = 0$ the capacitor currents to ground are:
\begin{center}
\renewcommand{\arraystretch}{1.3}
\begin{tabular}{cccc}
\toprule
Line & Cap & Voltage & Cap current (units of $j\omega C$) \\
\midrule
1 & $C$   & $+2/3\,\text{V}$ & $+2/3$ \\
2 & $C$   & $-1/3\,\text{V}$ & $-1/3$ \\
3 & $kC$  & $-1/3\,\text{V}$ & $-k/3$ \\
\bottomrule
\end{tabular}
\end{center}
Total ground current $= j\omega C\!\left(\tfrac{2}{3} - \tfrac{1}{3}
- \tfrac{k}{3}\right) = j\omega C\cdot\dfrac{1-k}{3}$.
\begin{itemize}
    \item $k = 1$: ground current $= 0$.  Perfect cancellation. \checkmark
    \item $k = 0.5$: ground current $= j\omega C/6$.
          Matches $I_{\mathrm{CM}} = j\omega C/6$. \checkmark
\end{itemize}
Cap 3 draws less current than it would if $k = 1$.  The
``shortfall'' escapes as $I_{\mathrm{CM}}$ — the conducted EMI
current a LISN would measure.
\end{physbox}

\begin{statebox}[Summary: $k = 1$ vs.\ $k = 0.5$ for $N = 3$]
\begin{center}
\renewcommand{\arraystretch}{1.5}
\begin{tabular}{lcc}
\toprule
\textbf{Current}
    & $k = 1$ (symmetric)
    & $k = 0.5$ (asymmetric) \\
\midrule
$I_{12}$ (DM pair 1--2)
    & $j\omega C/2$
    & $j\omega C/2$ \emph{(unchanged)} \\
$I_{23}$ (DM pair 2--3)
    & $0$
    & $-j\omega C/12$ \emph{(appeared!)} \\
$I_{\mathrm{CM}}$
    & $0$
    & $j\omega C/6$ \emph{(appeared!)} \\
$|I_{\mathrm{CM}}|$ at $C=1\,\text{nF}$, $1\,\text{MHz}$
    & $0\,\text{mA}$
    & $1.05\,\text{mA}$ \\
Modes coupled?
    & \textbf{No}
    & \textbf{Yes} \\
\bottomrule
\end{tabular}
\end{center}
\textbf{New insight beyond $N = 2$:} asymmetry on a single line
couples not just DM to CM, but also the two DM pairs to each other.
$I_{23} \neq 0$ even though $V_{23} = 0$ — the DM(1--2) oscillation,
filtered through the cap mismatch, induces a DM(2--3) current as a
side effect.  DM pair (1--2) itself is oblivious (its equation is
$k$-independent), yet it excites the rest of the system.
\end{statebox}

\begin{physbox}[The Big Picture]
The transformation matrices $\TiN$ and $\TvN$ have now done their
full job.  They converted the messy $N \times N$ mixed-mode
equation into a structured set where:
\begin{enumerate}
    \item $R$ and $L$ contribute only diagonal terms (via $\mathbf{M}$)
          — they never cause DM--CM coupling.
    \item The capacitor matrix $\mathbf{S}_N$ is the sole source of
          coupling, and only through the off-diagonal entries of
          $\mathbf{S}_{\mathrm{DCM}}$.
    \item The coupling magnitude is proportional to $(k-1)$:
          the further the circuit deviates from symmetry, the stronger
          the DM-to-CM noise injection.
\end{enumerate}
The next sections apply this framework to an actual SiC half-bridge
power module, where specific values of $k$ and the parasitic
capacitance $\Cp$ allow us to build a concrete CM equivalent circuit.
\end{physbox}

% \section{The Half-Bridge Testbed and the MCPM}
% TODO: §3 — MCPM structure, Fig. 2 walk, LISN explanation,
%            reference A, three-section split (Fig. 3)

% \section{The First CM Equivalent Model}
% TODO: §4 — LISN CM equiv., asymmetric cap CM eq. (Eq. 19),
%            reformulation (Eqs. 20-23), Fig. 4 walk, approx. (Eqs. 24-26),
%            Fig. 5 validation, Tables I-II

% \section{The Alternative CEM: Same Answer, Different Insight}
% TODO: §5 — two-section split (Fig. 6), alternative CEM (Fig. 7),
%            leakage vs. recycling comparison
